%% =====================================================================
%%  oficio.tex  --  MAESTRO del oficio (plantilla del despacho)
%% ---------------------------------------------------------------------
%%  Flujo de trabajo:
%%    1. Copia esta carpeta:  documentos/oficios/YYYY-MM-DD-titulo/
%%    2. Edita metadata.tex (numero, fecha, ...) y el cuerpo de abajo.
%%    3. Compila:  scripts/build.sh <carpeta>/oficio.tex
%%
%%  Selecciones disponibles (descomentar para cambiar el default):
%%    \usarmembrete{despacho}     % logo Camara de Diputados (moderno)
%%    \sinmarcaagua               % sin marca de agua
%%    \setmarcaaguaancho{10cm}    % marca de agua mas pequena
%%    \confolio                   % "Pag. 1 de N" (oficios largos)
%% =====================================================================
\documentclass{oficio}

%% =====================================================================
%%  metadata.tex  --  DATOS del documento (lo unico que se edita aqui)
%% ---------------------------------------------------------------------
%%  El estilo lo pone el framework; este archivo solo declara datos.
%%  Los defaults del despacho (autor, cargo, lema, siglas) ya vienen
%%  del nucleo: sobrescribir solo lo que cambie.
%% =====================================================================

\setmeta{numero}{02}
\setmeta{anio}{2026}
% \setmeta{siglas}{ABC/CR}          % default del nucleo
\setmeta{lugar}{Lima}
\setmeta{fecha}{30 de julio del 2026}   % comentar para usar la fecha de hoy

%% Nombre institucional del PDF (estandar del despacho):
%%    TIPO N°XXXX-AAAA - Sumilla.pdf   (se compone al compilar)
%% Descomentar y poner la descripcion breve / destinatario:
% \setmeta{sumilla}{DREA}
%% O bien un nombre totalmente a medida:
% \setmeta{nombrearchivo}{OFICIO N°0015-2026-2027 - DREA}


\begin{document}

\lugaryfecha
\numerodocumento

\destinatario{Señor:}{Xxxxxxxx}{}

\asunto{Invita al plenario ``Encuentro por la educación ayacuchana''}

\saludo

Es un honor saludarle cordialmente a nombre del \meta{cargo},
\meta{autor}, y expresar el reconocimiento a su
destacada labor a favor de nuestra comunidad educativa.

El motivo del presente documento es para invitarle de manera especial al
plenario ``Encuentro por la educación ayacuchana'', un espacio para sumar
esfuerzos, articular estrategias y coordinar acciones concretas que
contribuyan al desarrollo de la educación ayacuchana.

El encuentro se llevará a cabo de acuerdo con la siguiente programación:

\begin{itemize}
  \item Fecha: Sábado, 08 de agosto del 2026.
  \item Hora: 9:00 a.\,m.
  \item Lugar: Auditorio de la Casa del Maestro, Jr.~2 de Mayo N.º~235.
\end{itemize}

Conociendo su constante compromiso con nuestra educación y el desarrollo
social, estamos seguros de que su participación y valiosos aportes
enriquecerán significativamente las decisiones que se adopten en este
espacio.

Sin otro particular por el momento, me despido reiterándole los
sentimientos de mi especial consideración y estima.

\atentamente
%% \vfill ancla la firma al FONDO de la pagina: misma altura en todos
%% los oficios, sin importar el largo del cuerpo. El argumento de
%% \bloquefirma es el minimo garantizado sobre el nombre (ahi va la
%% firma digital).
\vfill
\bloquefirma[2.0cm]

\end{document}
